% SPDX-License-Identifier: CC-BY-4.0
\documentclass[11pt,letterpaper]{article}

\usepackage[T1]{fontenc}
\usepackage{amsmath,amssymb}
\usepackage{tikz}
\usepackage{pgfplots}
\pgfplotsset{compat=1.18}
\usepackage{algorithm}
\usepackage{algpseudocode}
\usepackage{booktabs}
\usepackage{enumitem}
\usepackage[hidelinks]{hyperref}

\newcommand{\N}{\mathbb{N}}
\newcommand{\Z}{\mathbb{Z}}
\newcommand{\Q}{\mathbb{Q}}
\newcommand{\R}{\mathbb{R}}
\newcommand{\C}{\mathbb{C}}
\newcommand{\Map}[3]{#1 \colon #2 \to #3}
\newcommand{\Mod}[1]{#1\text{-}\mathrm{Mod}}
\newcommand{\GroupHom}{\mathrm{Hom}_{\mathbf{Grp}}}
\newcommand{\RingHom}{\mathrm{Hom}_{\mathbf{Ring}}}

\newtheorem{theorem}{Theorem}
\newtheorem{lemma}[theorem]{Lemma}
\newtheorem{corollary}[theorem]{Corollary}
\newtheorem{definition}[theorem]{Definition}
\newtheorem{remark}[theorem]{Remark}
\newenvironment{proof}{\noindent\textit{Proof.}}{\hfill\(\square\)}

\title{Generalized Collatz Library}
\author{}
\date{\today}

\begin{document}
\maketitle
\section*{Affine escape outside a window}

\begin{abstract}
For the positive raw map that divides by \(D\) on divisible inputs and applies
\(n\mapsto qn+c\) otherwise, a simple greatest-common-divisor obstruction
completely determines a large family of orbits. If
\(\gcd(q,D)\nmid c\), every positive seed has a finite initial run of divisions
followed by an affine tail. We prove the exact time-indexed formula for that
tail and deduce escape to infinity for every positive seed. The argument works
for composite \(D\), uses no probabilistic branch model, and states explicitly
what remains unresolved when the obstruction fails.
\end{abstract}

\section{Map and hypotheses}

Fix integers
\[
q\ge2,\qquad c\ge1,\qquad D\ge2.
\]
The positive raw map is
\[
F(n)=
\begin{cases}
n/D,&D\mid n,\\
qn+c,&D\nmid n,
\end{cases}
\qquad n\in\N_{>0}.
\]
Both branches preserve positivity. Put
\[
g=\gcd(q,D)
\]
and impose the arithmetic obstruction
\[
g\nmid c.                                                    \tag{1}
\]
Only elementary divisibility properties of the greatest common divisor are
used; standard background may be found in Hardy and Wright
\cite{HardyWright08}.

For \(j\in\N\), define the geometric recurrence
\[
G_0(q)=0,
\qquad
G_{j+1}(q)=qG_j(q)+1.                                         \tag{2}
\]
Induction in (2) gives
\[
G_j(q)=1+q+\cdots+q^{j-1}=\frac{q^j-1}{q-1}                  \tag{3}
\]
for \(j\ge1\), with the empty sum interpreted as zero.

\section{The division barrier}

\begin{lemma}[Gcd division barrier]
\label{lem:barrier}
Under (1), for every positive integer \(y\),
\[
D\nmid qy+c.
\]
\end{lemma}

\begin{proof}
Suppose instead that \(D\mid qy+c\). Since \(g\mid D\), transitivity of
divisibility gives \(g\mid qy+c\). Also \(g\mid q\), hence \(g\mid qy\).
Subtracting these two divisible integers gives \(g\mid c\), contrary to
(1). Therefore \(D\nmid qy+c\).
\end{proof}

\begin{remark}
The lemma is deterministic. It does not estimate how often divisions occur;
it proves that after any affine instruction no later division can occur. The
conclusion uses the actual divisor \(D\), not a prime factor or a randomized
residue model.
\end{remark}

\section{Canonical initial divisions}

\begin{lemma}[Whole-base factorization]
\label{lem:factorization}
For every \(n\in\N_{>0}\), there is a unique pair \((k,x)\) with
\[
k\in\N,
\qquad
x\in\N_{>0},
\qquad
n=D^k x,
\qquad
D\nmid x.                                                     \tag{4}
\]
For this pair,
\[
F^i(n)=D^{k-i}x\quad(0\le i\le k),                           \tag{5}
\]
and in particular \(F^k(n)=x\).
\end{lemma}

\begin{proof}
Consider the set
\[
S=\{r\in\N:D^r\mid n\}.
\]
It is nonempty because \(0\in S\). If \(r\in S\), positivity gives
\(D^r\le n\). Since \(D\ge2\), this bounds \(r\), so \(S\) has a largest
element \(k\). Write \(n=D^k x\). The quotient \(x\) is positive. If
\(D\mid x\), then \(D^{k+1}\mid n\), contradicting maximality. Thus
\(D\nmid x\).

For uniqueness, suppose also \(n=D^\ell y\) with \(D\nmid y\). If
\(\ell<k\), then \(y=D^{k-\ell}x\), so \(D\mid y\), a contradiction. The
same argument with the pairs reversed excludes \(k<\ell\). Hence
\(k=\ell\) and then \(x=y\).

For (5), induction on \(i\) suffices. The case \(i=0\) is (4). If \(i<k\)
and \(F^i(n)=D^{k-i}x\), then the exponent \(k-i\) is positive, so this
state is divisible by \(D\). The raw map therefore selects its division
branch and
\[
F^{i+1}(n)=D^{k-i-1}x.
\]
This proves (5) through \(i=k\).
\end{proof}

The exponent \(k\) in Lemma~\ref{lem:factorization} is a whole-base order.
No primality or prime-power hypothesis on \(D\) is present. For example, if
\(D=6\), then \(k\) counts powers of \(6\), not the separate \(2\)-adic or
\(3\)-adic valuations.

\section{Exact affine tail and escape}

\begin{theorem}[All-seed affine escape]
\label{thm:escape}
Assume (1), let \(n\in\N_{>0}\), and let \((k,x)\) be the unique pair from
Lemma~\ref{lem:factorization}. Then for every \(j\in\N\),
\begin{align}
F^{k+j}(n)&=q^j x+cG_j(q),                                   \tag{6}\\
D&\nmid F^{k+j}(n).                                           \tag{7}
\end{align}
Consequently,
\[
F^{k+j}(n)\longrightarrow\infty
\qquad\text{as }j\longrightarrow\infty.                    \tag{8}
\]
Thus every positive seed escapes to infinity.
\end{theorem}

\begin{proof}
By Lemma~\ref{lem:factorization}, \(F^k(n)=x\) and \(D\nmid x\). This proves
(6) and (7) at \(j=0\), because \(G_0(q)=0\).

Assume (6) and (7) at an index \(j\), and write
\[
y=F^{k+j}(n)=q^j x+cG_j(q).
\]
Since \(D\nmid y\), the next raw instruction is affine. Therefore
\begin{align*}
F^{k+j+1}(n)
 &=qy+c\\
 &=q\bigl(q^j x+cG_j(q)\bigr)+c\\
 &=q^{j+1}x+c\bigl(qG_j(q)+1\bigr)\\
 &=q^{j+1}x+cG_{j+1}(q).
\end{align*}
This is (6) at \(j+1\). Lemma~\ref{lem:barrier}, applied to the positive
state \(y\), shows that this new affine output is not divisible by \(D\).
Hence (7) also holds at \(j+1\), completing the simultaneous induction.

All terms in (6) are nonnegative and \(x\ge1\), so
\[
F^{k+j}(n)\ge q^j x\ge q^j.
\]
Because \(q\ge2\), the sequence \(q^j\) tends to infinity. This proves (8).
The seed \(n\) was arbitrary, so the conclusion holds for every positive
seed.
\end{proof}

\section{Example and exact recovery}

Take
\[
(q,c,D)=(6,1,4),
\qquad
n=64.
\]
Here \(g=\gcd(6,4)=2\nmid1\). The canonical factorization is
\(64=4^3\cdot1\), so \(k=3\) and \(x=1\). The initial prefix is
\[
64\longmapsto16\longmapsto4\longmapsto1.
\]
Theorem~\ref{thm:escape} now gives
\[
F^{3+j}(64)=6^j+G_j(6).
\]
For \(j=0,1,2,3,4\), this produces
\[
1,
\quad 7,
\quad 43,
\quad 259,
\quad 1555.
\]
Every number in this tail is congruent to either \(1\) or \(3\) modulo
\(4\), as guaranteed by the theorem.

The recovery statement is exact. Given \(n\), the factorization (4)
reconstructs every state before time \(k\) via (5), while (6) reconstructs
every state from time \(k\) onward. The branch guards are also recovered:
times before \(k\) use division, and all later times use the affine branch.
No subsequence, accelerated clock, or discarded residue class occurs.

\section{Reduction certificate and boundaries}

The original obligation in the parameter region (1) is to classify every
positive raw orbit. The replacement obligations are finite extraction of the
maximal initial power \(D^k\) and evaluation of (6). Lemma
\ref{lem:factorization} and Theorem~\ref{thm:escape} prove recovery in both
directions of time indexing: together they cover every raw time and preserve
every branch condition.

A concrete resource strictly decreases. Before using (1), an orbit may appear
to have unboundedly many unresolved future choices between division and the
affine branch. After the finite prefix, Lemma~\ref{lem:barrier} reduces the
number of unresolved future branch choices to zero. The remaining recurrence
is explicit and its lower bound proves (8). The argument therefore closes the
all-seed obligation in this region rather than transferring it to another
unknown dynamical system.

The hypothesis (1) is sufficient, not necessary. When \(g\mid c\), the proof
of Lemma~\ref{lem:barrier} yields no contradiction, and this article makes no
claim about the resulting cell. In particular, the classical cell
\((q,c,D)=(3,1,2)\) cannot satisfy (1), because its gcd is one. The theorem
also does not cover signed states, signed coefficients, zero seeds, or
\(q<2\). These are domain boundaries, not conclusions about those omitted
systems.

\bibliographystyle{amsalpha}
\bibliography{AffineEscape}

\end{document}
